\section{Results}
\label{sec:results}

\subsection{Dataset and protocol}
\label{sec:results/datasets}

We measure on the \textbf{CESNET-QUICEXT-25} XS subset, a publicly released
year-long QUIC traffic corpus with extended flow attributes
\citep{cesnet2025cesnet}. From it we take the paired-design subset: captures
carrying service labels, restricted to the services and flows on which the
occupancy gate is defined.

The \textbf{990-pair denominator} comes from a flow-count threshold applied
before any bracket is computed. Of the 53 service labels present, 45 carry at
least 2{,}000 flows in the retained period, and those 45 form
$\binom{45}{2} = 990$ unordered pairs, all surviving the hygiene filters. Flow
segmentation and service selection follow a pre-registered rule
(\Cref{sec:appendix/preprocessing}), and the 2{,}000-flow threshold was fixed before the gate was
measured; the subset pass is provisional rather than a full-corpus confirmation.

The masking regime is verified rather than assumed: the padding sweep reports a
single zero-valued bin across the sampled tail slots, 990 usable pairs at
$k{=}5$, no services lost to the prefix-length filter, and identical masked and
unmasked symbol counts at every rung. A bracket over a channel that silently
included padding or address structure would be a different channel than the one
declared.

The headline $336/990$ is computed under the full-$n$ convention: every eligible
flow contributes, with no per-service cap. Both corpora were measured under both
conventions; the primary corpus's capped variant gives $335/990$ rather than
$336/990$, with pair $(21,23)$ the sole difference
($\epsilon_{\rm lo}^{\rm capped} = 0.0498$ falls below the $0.05$ gate while
$\epsilon_{\rm lo}^{\rm full\text{-}n} = 0.0588$ clears it). The
$20{,}000$-flow cap binds on 814 of the 990 primary-corpus pairs; of those,
capping raises $\epsilon_{\rm lo}$ on 38 (a 4.7\% reversal rate within the
cap-bound subset, 3.8\% of all pairs). The aggregate direction holds: 776 pairs
see $\epsilon_{\rm lo}$ decrease under capping against 38 where it increases.

\subsection{The two-sided bracket}
\label{sec:results/bracket}

\Cref{tab:bracket} (in \S A.4) is the paper's primary result. At the coarsest rung
($k{=}5$, $M{=}1$), 336 of the 990 pairs satisfy the joint gate of
\eqref{eq:gate}, a fraction of $33.9\%$. The two sides behave differently, and
that difference is the point: the upper side admits 340 pairs ($0.3434$) against
the lower side's 986 ($0.9960$), and the four-pair gap is exactly the set
passing the upper side and failing the lower. At the coarsest rung, where the joint gate is defined, the upper
side is therefore the scarce object: it is what fails relative to the lower
side, and it is what makes the certified set smaller than either side alone. At
finer resolution the lower side is not what binds: on the m-free route at
$(M{=}2, k{=}5)$ every pair clearing the upper bound also clears the lower gate,
so both counts are $127$, and on the m-free route at $(M{=}1, k{=}10)$ both are
$98$. The coarsest cell is the exception, the lower side
narrowing it by four pairs ($340\to336$). The invariant $\epslo \le \epshi$ has
zero violations on $990/990$ rows at every cell, each row's m-free upper bound
recorded. The smallest lower bound among the $127$ is $0.0755$, $0.0255$ above the
gate, against a largest measured convention effect of $0.0072$: $3.5\times$
clearance, so the identity is robust to that convention. The upper-side column of \Cref{tab:bracket} is what falls to
zero as the alphabet is refined; \Cref{sec:results/horizon} shows that the cause
is alphabet sparsity in the plug-in estimate at fine resolution rather than the
alphabet-priced radius.

The headline fraction is a stratum statement, not a property of traffic in
general: it says that 33.9\% of pairs on this corpus, at this channel, admit a
two-sided interval whose upper side is finite-sample and whose lower side clears
the gate. It does not say the remaining 66.1\% are indistinguishable: for most of
them no certificate \emph{exists} here, which \Cref{sec:results/horizon} shows is
a different and stronger statement.

\paragraph{Uncertainty on the headline.}
The 990 pairs are not independent observations: they form a complete graph over
45 services, so pairs sharing a service are correlated and a pair-level interval
would be too narrow. We therefore resample \emph{services} and report
service-level cluster-bootstrap intervals. The joint gate of $336/990 = 33.9\%$
carries a $95\%$ interval of $[26.5\%, 43.1\%]$; the upper side
$[26.8\%, 43.4\%]$; and the lower side $[99.1\%, 100\%]$. The gate interval is
wide, which is the honest statement about a corpus of 45 services; the lower
side's is tight, the same observation the gate makes, that most pairs
are separable here and the certificate is limited by the upper side.

Saturation is zero at this rung: no pair's plug-in estimate $\widehat{tv}$
reaches $0.99$, so the informative fraction is the whole corpus and the
certificate is not being purchased by discarding saturated pairs.
\Cref{tab:endpoints} in the appendix gives both endpoints' quantiles.

The same measurement on an independent corpus yields a joint gate of
$523/3321$ ($15.75\%$), an upper-side fraction of $0.1608$ and a lower-side
fraction of $0.9967$. The two results are analysed together in
\Cref{sec:discussion/robustness}, where the sample-size stratification behind the
difference, the conservative direction of its mixed-side sample sizes, and a
service-cluster bootstrap of the gap are set out.

A second certified object is computed on the same flows over whole-flow rectified
statistics rather than the $k$-prefix symbol law. At $M{=}1$ its upper side admits
292 of 990 pairs and its joint gate 289, against the prefix law's 340 and 336.
\Cref{sec:discussion/objects} shows the proximity of their informative fractions
to be cancellation, not agreement: we neither rank them nor treat either as a
check on the other.

\subsection{The certificate horizon}
\label{sec:results/horizon}

Increasing the alphabet resolution makes the channel more informative and the
certificate less available, and both effects are measured. The median $\tv$ over
pairs rises with $M$, from $0.5805$ at $M{=}1$ through $0.8833$, $0.9248$, and
$0.9462$ to $0.9673$ at $M{=}16$.

The rise is tested with paired increments rather than by comparing marginal
intervals, since the same service draws are reused at every rung and the
marginals overlap at the finer resolutions. A paired increment at a step is the
median of the paired differences
$\Delta_i = \widehat{tv}(M,\ \text{pair}_i) - \widehat{tv}(M-1,\ \text{pair}_i)$
taken over pairs, \emph{not} the difference of the two rungs' medians; the two
statistics differ because the median is not additive, and at the first step they
give $+0.2931$ (median of differences, what we report) against $+0.3028$
(difference of the two rungs' median divergences, $0.8833 - 0.5805$ from the
first line above). The paired form is the one that controls for pair identity
and carries the interval.
The paired increment excludes zero at every step: $+0.2931$ $[+0.2580,+0.3320]$ from $M{=}1$
to $2$; $+0.0438$ $[+0.0229,+0.0708]$ to $4$; $+0.0199$ $[+0.0040,+0.0429]$ to
$8$; and $+0.0131$ $[+0.0003,+0.0403]$ to $16$. The paired increment excludes
zero at every rung and shrinks by over an order of magnitude across the ladder,
though the last step is the thinnest, its lower bound
essentially on zero, so we describe the tail as still rising but flattening
rather than asserting it as firmly as the first step. One confound must be named: the increments establish that the medians rise with
significance, but do not by themselves separate genuine resolution from sparse
inflation at fixed $n$, since a finer alphabet spreads the same flows over more
symbols. The saturating column is the tell (\Cref{tab:bracket} / \Cref{fig:horizon}a). The $M{=}1$ median's marginal interval is
$[0.5079,0.6459]$.

A finer alphabet is associated with higher median divergence, but the confound
named above leaves the causal claim open. What falls is the fraction of pairs for
which an upper bound
below the gate threshold still exists: $340/990$ at $M{=}1$, $47/990$ at
$M{=}2$, $3/990$ at $M{=}4$, and zero at $M{=}8$ and $M{=}16$
(\Cref{fig:horizon}a).

We first attributed this to the radius, and that attribution was wrong.
The certified floor is $(1-\epshi)/2$ with $\epshi=\widehat{tv}+r_a+r_b$, and
$r_i$ in \eqref{eq:upper} does grow with $m=(2M)^k$. But the estimator has
already saturated before any radius is added: at $k{=}30$, $597$ of $820$ pairs
carry $\widehat{tv}\ge 0.99$ with median $0.9979$, leaving no room under the
gate threshold whatever the radius. The cause is the sparsity-driven upward bias
of $\widehat{tv}$ at fine resolution, not the alphabet-priced radius.

That distinction matters because the alphabet price is avoidable. Our
\eqref{eq:upper} bounds the $\ell_1$ norm, whose supremum runs over all $2^m$
subsets, and $\sqrt{m/n}$ is that route's cost. Bounding instead the single
population-optimal set $A^\ast=\{x:P(x)>Q(x)\}$, which depends only on the true
$P,Q$, is a two-event statement rather than a $2^m$-event one, so one-sided
Hoeffding applies directly to $P(A^\ast)$ and $Q(A^\ast)$ and no $m$ appears:
\begin{equation}
  \tv(P,Q) \;\le\; \widehat{tv}
  + \sqrt{\tfrac{\ln(2/\delta)}{2n_a}} + \sqrt{\tfrac{\ln(2/\delta)}{2n_b}} .
  \label{eq:mfree}
\end{equation}
This is a route choice, not a new theorem: both bounds are valid and certify
the \emph{same} object, an interval on $\tv(P,Q)$ for the pair, differing only in
derivation. The two-event deviation argument is the Le Cam mechanism already
used above \citep{tsybakov2009introduction}, and avoiding the alphabet term this
way is the theme of the closeness-testing line
\citep{chan2013closeness,diakonikolas2020highprob} rather than a device particular
to this paper. At our per-law level $\delta=\alaw=0.025$ the fluctuation terms
coincide exactly, $\tfrac12\sqrt{2\ln(2/\alaw)/n}=\sqrt{\ln(2/\alaw)/(2n)}$, so
the two forms differ only in $\sqrt{m/n}$ and no constant convention needs
reconciling.

Both are reported. On the same per-pair data \eqref{eq:mfree} is $2.9\times$
tighter at the coarsest rung ($M{=}1$, $k{=}5$, $n{=}2791$; per-side radius
$0.028018$ against $0.081557$), and the nominal coverage measured by the calibration design of
\Cref{tab:coverage} applies to the alphabet-priced upper side. Recomputed over the full
$(M,k)$ grid (\Cref{tab:grid}), the two bounds show two regimes rather than one. At coarse resolution the
alphabet price is what binds, and removing it buys real headroom: at the coarsest
rung ($M{=}1$) the frontier reaches $k{=}15$ with $18$ pairs certifiable where
the published bound reports zero, and the m-free route also clears $4$ pairs at
$M{=}2$ and $1$ at $M{=}4$ at that $k$. The advantage does not generalize to the finer
prefix lengths: it is present at every $M$ for $k{=}5$ and $k{=}10$, is confined
to $M{=}1,2,4$ at $k{=}15$, and is absent at $k{=}30$, so the extension is a
property of the coarse rungs rather than of the frontier as a whole. At fine resolution
sparsity reasserts the frontier regardless of the radius's form, though the
transition is gradual rather than a threshold: at $k{=}10$, saturation is
already $59\%$ at $M{=}4$ while the m-free route still certifies $10$ pairs.
When $m \gg n$ most symbol cells hold $0$ or $1$ observation, so $\widehat{tv}$
approaches $1$ mechanically: estimator bias, not a property of the traffic. A
simulation at $n{=}2791$ confirms it: at $m{=}32$ the estimator is unbiased,
while at $m{=}10^6$ its median exceeds $0.99$ even at true TV $0.2$. At
$k{=}15$, $M{=}8$, $712$ of $946$
pairs ($75\%$) already carry $\widehat{tv}\ge0.99$ with median $1.0000$, so no
bound fits under the gate whatever its construction, and the saturated fraction
climbs monotonically with $M$ at every $k$. The horizon is therefore not dissolved by dropping the alphabet
price; the price binds at the coarse end, and saturation reasserts the boundary
at the fine end, which is a boundary of the object rather than a defect of the
construction. We keep
\eqref{eq:upper} primary because the paper's counts and validations were computed
on it, and report \eqref{eq:mfree} alongside rather than substituting, so the
horizon claim rests on the comparison rather than a swap. Every cell of
\Cref{tab:grid} comes from the released per-pair ledgers, which carry all five
$M$-rungs for the same pairs at every $k$; we report the grid as measured and
interpolate nothing.

Where the plug-in certificate is lost, the large-$k$ side is not empty. A
divisibility-based lower bound, evaluated on the rung set
$k \in \{5,10,15,20,25,30\}$ (the chain-rule experiment's grid, which extends the
bracket's), is nontrivial on every pair at $k{=}15,20,25,30$, so the pivot from a
two-sided bracket to a one-sided statement is available precisely where the
certificate fails. We report this bound as conditional on its realizability
assumption and never as distribution-free.

\subsection{Detector validation}
\label{sec:results/detector_validation}

A bracket is useful only if real detectors respect it. We test this with a
pre-registered one-sided binomial comparison on the informative stratum
($\epshi < 0.5$) against the predicted floor $(1-\epshi)/2$. The stratum is populated at $k{=}5$ (340 informative, 540 secondary, 110
excluded), summing to the 990-pair denominator of \Cref{tab:bracket}. At $k{=}15$ and $k{=}30$ the stratum is empty, so those
rungs contribute no validation evidence and we say so rather than reporting a
vacuous pass.

\paragraph{Panel and protocol.}
The pre-registration was written before the first measurement, with seed 12 and
flow- and time-disjoint splits; the registration timestamp precedes the
first-measurement timestamp. Three detectors are executed and one is not.
\emph{LightGBM} is trained on the $k{=}5$ (direction, size) feature vector, with
100 boosting rounds, learning rate $0.1$, and 31 leaves, on CPU. The detector's
features are the raw direction bits and log-transformed packet sizes at each of
the $k{=}5$ prefix positions (10 dimensions), which is a strictly finer
observable than the $m{=}(2M)^k{=}32$-cell quantized alphabet on which the
certificate's radius is computed. A detector that fails to breach the certified
floor is therefore evidence about a more informative observable than the one the
bracket covers. The resolution substitution runs in the direction of an easier
test: at finer resolution both the plug-in divergence and the confidence radius
increase, which lowers the predicted floor and makes a violation easier to
observe, so a non-breach on the finer object is the stronger statement. A breach
on the finer channel would not, however, falsify the coarser certificate; this
validation is a consistency check, not a falsifiable test. The closed-set /
open-set substitution declared below runs the other way, narrowing the citable
statement to the surrogate actually measured.
\emph{Prototype} is a nearest-centre rejector in the style of prototype-radius
open-set methods \citep{qu2026sepspace}: one mean centre per service, with a
per-service acceptance radius set at a distance quantile $0.95$ of the training
points around its own centre, so a test flow is accepted when its distance to
the nearest centre falls inside that centre's radius.
\emph{RoNeTC} follows the uncertainty-quantified rejector of
\citet{wang2025reliable}: an MLP with two hidden layers of width
$\min(128, 2d)$, ReLU activations and dropout $0.20$, trained with Adam at
$10^{-3}$ for $5$ training epochs and scored by Monte-Carlo dropout sampling
($8$ passes), with rejection driven by predictive uncertainty.
A fourth detector, a Pcap-Encoder adapter, was not part of the evaluated
panel; it contributes no cells and is not counted as a pass.

\paragraph{Multiplicity and power.}
The nine cells are three detectors crossed with three holdout levels, and the
same 340 informative pairs are reused across all nine. We report each cell's
raw one-sided $p$-value against its own null and do not combine cells; the
strongest cell is $p = 0.9765$, so no correction for multiplicity could change
the verdict, and we make no family-wise claim beyond that. The binomial treats
the 340 pairs as independent draws, which they are not: they form a graph over
the same services, with a maximum pair-degree of 27, so the effective
independent sample size lies somewhere between 43 and 340 (a heuristic bound,
used conservatively). The $340$
informative pairs span $43$ of the $45$ services, under the full-$n$ convention
of \Cref{sec:results/datasets} (the other $2$ carry no informative pair). The reported
$p$-values are therefore anti-conservative. That does not affect the
no-rejection conclusion, which holds by the widest margin at the LightGBM
holdout-$0.2$ cell ($p = 0.9765$ in \Cref{tab:detectors}), but
it does mean the test has less power than the nominal figure suggests.

Power is the more serious limitation. The test's null rate is
$\ptest = 0.05$, distinct from the confidence level $\alpha$ used elsewhere, and
8 violations in 340 pairs gives $p = 0.9954$, so the design has little power
against a modest excess. The observed rates ($0.6\%$--$2.9\%$) sit below both
null rates. A threshold of $0.025$ is implied by the joint level rather than
pre-registered; under it the LightGBM holdout-$0.2$ cell, the most adverse of the
nine, gives $p = 0.346$ on its $10$ violations.
Neither rejects, so reading ``not falsified'' as evidence of a tight floor would
be a mistake: the test supports only that no gross violation is present, and
detecting a two- or three-point excess would need more informative pairs than
this channel supplies.

\Cref{tab:detectors} reports the outcome. No cell rejects the null: observed
violation rates run from $0.6\%$ to $2.9\%$ against a null rate of $5\%$, all
nine cells are below the null, so no gross violation of the floor is present. Two
qualifications travel with the table, and the first is a scope restriction
rather than a caveat. The compared quantity is a pair-level macro-average
surrogate: it is the mean of the two constituent services' \emph{closed-set}
classification-error rates, whereas the certified floor of \eqref{eq:floor} is an
\emph{open-set} error bound. The test therefore does not measure the open-set
quantity the floor bounds; the citable statement is the narrower one, that
per-pair macro-average closed-set error stays above the floor on the informative
stratum.

 A fourth detector in the intended panel is not
reported: it did not run, its cells carry no information, and it is not counted
as a pass.

\emph{A consistency check on the certificate's own object.} A logistic
regression on the same $32$-symbol one-hot representation reaches a median
balanced error of $0.366$ against a per-pair floor $(1-\epshi)/2$ (median
$0.198$), clearing its own floor on every pair ($0/990$); its implied $\tv$
(median $0.267$) is weaker than the certified lower bound (capped-convention
median $0.535$), so
the certificate is consistent with, but not tight against, the classification
risk.

\subsection{The standard estimator cannot answer the growth question}
\label{sec:results/estimator}

The remaining question is how divergence accumulates as the prefix grows: with
independent positions $D_k$ would grow linearly, and sublinear growth would mean
later packet positions add less than the first. The natural estimator is a
chain-rule plug-in over conditional laws, and we admitted it to real data only
after it passed a known-truth control whose per-position divergence is known
analytically, exactly $0.3389191441548814$, so its joint divergence grows
linearly by construction.

It did not pass, and that failure is the finding. On that analytically linear
source the estimator returned sublinear growth in both arms, including a twin
built on the real occupancy profile. Its estimate-to-truth ratio is not monotone:
it peaks slightly above $1.0$ mid-range then falls to $0.6769$ on the plain arm
and $0.8820$ on the twin at $k{=}30$, deflating divergence where the prefix is
longest (\Cref{tab:estimator} and \Cref{fig:estimator}a, both in \Cref{sec:appendix/limitations}). The mechanism is
measured, not asserted: occupancy collapses with $k$, the fraction of contexts
with fewer than five transitions rising from $0.000$ at $k{=}5$ to $0.998$ at
$k{=}30$, and the median transitions per context falling from $200{,}000$ to $1$
(\Cref{fig:estimator}b,c). A per-position smoothing bias accumulates where
contexts are sparse, biasing the fitted exponent downward.

Two independent repairs were attempted and both failed their known-truth gate.
Replacing the estimator family with a Laplace-corrected, jackknife-debiased
estimator failed on both arms. Conditioning on exposure, which restricts
attention to contexts with at least $T$ transitions, never passed at any
$T \in \{5,10,20,50,100\}$ on the plain arm, and the excluded mass needed to make
contexts well-occupied grows to $0.608$ at $k{=}30$. That last number is the
finding's teeth: the repair that would work costs more than half the data at
the upper end of the prefix range. The obstacle is structural, not a matter of
tuning, and the honest statement is that the growth question is open. A
practitioner who needs a growth-rate answer on traffic of this sparsity needs a
different measurement design, and this section says so with numbers instead of
asserting it.

One further result is reported only to show the question is non-vacuous. A
nested-model comparison of a first-order Markov law against an i.i.d. null
favored the Markov model for 45 of 45 services at every tested prefix length.
That vote is not evidence: a heterogeneous arm, in which each flow draws its own
rate and then emits symbols conditionally independently at that rate, reproduces
the same vote at $k \ge 8$. That arm is homogeneous \emph{within} a flow but not
across flows, and its symbols are not independent: two draws from one flow share
a rate, so their covariance is $\mathrm{Var}(\theta) > 0$. It is therefore an
alternative mechanism for dependence rather than a null that excludes it, and the
vote cannot separate the two explanations. The vote is consequently not
diagnostic, and we report it as a measurement artifact
rather than as support for dependence. The magnitude of the real-data statistic
does exceed what the null reproduces, so neither ``dependence is real'' nor
``heterogeneity explains it'' is the correct summary; what is established is
that the vote cannot settle the question.

